\begin{definition} [Saddle Point]
    Consider a function $f:\set{w} \times \set{z} \to \field{R}$.
    We refer to a pair $(w^*, z^*)$ as a \emph{saddle point} for $f$, if we have
    \[
        f(w^*, z) \leq f(w^*, z^*) \leq f(w, z^*),
    \]
    for all \inSet{w} and \inSet{z}.
\end{definition}

\begin{theorem}
    \label{th:saddle_is_central}
    Let $X$ be a discrete random variable, and \setQ be a finite set of distributions for $X$.
    Consider the following function:
    \begin{equation*}
        f(P, R)= \fExp,
    \end{equation*}
    where $Q_i \in \setQ$, and $P$, $R$ are probability distributions.

    If \saddle is a saddle point of $f$, then \saddP is a central distribution of \setQ for $X$.
\end{theorem}

\begin{proof}
    Since $X$ is a discrete random variable, \setQ has a central distribution for $X$.
    Let \cDist be a central distribution of \setQ, and \saddle be a saddle point for $f$.
    Clearly, the point $(\cDist, \saddR)$, is in the domain of $f$, and
    since \saddle is a saddle point for $f$, we have
    \begin{equation}
        \label{eq:th_proof_saddle_lq_dx}
        f\saddle \leq f(\cDist, \saddR) \cdot
    \end{equation}
    Because \saddR is a probability distribution, the definition of $f$ implies:
    \[
        f(\cDist, \saddR) \leq \supD[\cDist] \cdot
    \]
    On the other hand,
    \[
        f\saddle = \max_{R} f(\saddP, R) = \supD[\saddP] \cdot
    \]
    Thus, Inequality~\ref{eq:th_proof_saddle_lq_dx} implies
    \[
        \supD[\saddP] \leq \supD[\cDist] \cdot
    \]
    Consequently, $\saddP$ is also a central distribution of \setQ for $X$.
\end{proof}

\begin{lemma}
    \label{lm:saddle_conditions}
    Let $f$ be the function defined in theorem~\ref{th:saddle_is_central}.
    The point \saddle is a saddle point for $f$, if and only if there are constants $\lambda$, $\eta$ and a
    function $\mu$, such that for every $i$, the following conditions hold:
    \begin{gather}
        \label{eq:kkt_lambda}
        \klExp[Q_i][\saddP][x] - \mu(i) = \lambda, \\
        \label{eq:kkt_w_avg}
        \saddP(x) = \sum_i \saddR(i) Q_i(x) \\
        \label{eq:kkt_r_sums_one}
        \sum_i \saddR(i) = 1, \\
        \label{eq:saddle_lm_r_geq_0}
        \saddR(i) \geq 0, \\
        \nonumber
        \mu(i) \geq 0, \\
        \label{eq:saddle_lm_rmu_is_zero}
        \mu(i) \saddR(i) = 0 \cdot
    \end{gather}
\end{lemma}

\begin{proof}
    Similar to $f$, we define $\hat{f}$ as follows:
    \begin{equation*}
        \hat{f}(P, R) = \fExp,
    \end{equation*}
    where $P$, $R$ are arbitrary real functions instead of probability distributions.
    This enables us to formulate the saddle point of $f$ as the solution to a constrained optimization problem.

    Let \saddle be a simultaneous solution to the following two optimization problems:

    \begin{equation}
        \begin{aligned}
            &\text{minimize} & &\hat{f}(P,R) & &\text{w.r.t.} \ P \\
            &\text{subject to} & &\sum_{x} P(x) = 1, \\
            \\
            &\text{maximize} & &\hat{f}(P,R) & &\text{w.r.t.} \ R \\
            &\text{subject to} & &R(i) \geq 0, & &\text{for all} \ i  \\
            & & &\sum_{i} R(i) = 1 &
        \end{aligned}\label{eq:twin_optimization}
    \end{equation}
    Clearly, if $\saddle$ exists, it will be a saddle point for $f$.

    Both of these optimization problems are convex~\cite{0521833787} and differentiable.
    To establish this, consider that all the constraint functions are affine, and therefore are convex and
    differentiable.

    Additionally, $\hat{f}$ is differentiable function with respect to both $R$ and $P$.
    For any fixed $P$, $\hat{f}(P,\, \cdot \,)$ is an affine function and is concave.

    On the other hand, for a fixed solution to the second optimization problem denoted by $\hat{R}$, we can write
    \begin{equation}
        \label{eq:simplified_f}
        \hat{f}(P,\hat{R}) = c - \sum_{i,x} \hat{R}(i) Q_i(x) \ln P(x),
    \end{equation}
    where $c$ is not a function of $P$.
    Since $\hat{R}$ is a solution to the second optimization problem, it satisfies the constraints
    of the second problem and therefore is a probability distribution.
    Hence, for all $i$, $x$ we have $\hat{R}(i)Q_i(x) \geq 0$, and $\hat{f}(P,\, \cdot \,)$  is
    the negative of a linear combination of concave $\ln P(x)$ functions with nonnegative coefficients.
    This function is convex.

    Thus, both optimization problems are convex and differentiable, with affine constraints.
    The constraints simply restrict the feasible set to the set of probability distributions, and it is always possible
    to find feasible points that meet all the constraints.
    Therefore, Slater’s condition is satisfied, and the KKT conditions provide necessary
    and sufficient conditions for optimality~\cite{0521833787}.
    In other words, If a point $\saddle$ satisfies the KKT conditions for both problems, it will be a saddle
    point for $f$, and vice versa.

    The KKT conditions are as follows for every $x$, $i$ and constants $\lambda$, $\eta$:
    \begin{gather}
        \label{eq:kkt_proof_R_is_pxIe}
        \sum_i \saddR(i) Q_i(x) = \eta \saddP(x) \\
        \label{eq:kkt_proof_p_sums_one}
        \sum_{x} \saddP(x) = 1, \\
        \nonumber
        \klExp[Q_i][\saddP][x] - \mu(i) = \lambda, \\
        \label{eq:kkt_proof_r_sums_one}
        \sum_i \saddR(i) = 1, \\
        \nonumber
        \saddR(i) \geq 0, \\
        \nonumber
        \mu(i) \geq 0, \\
        \nonumber
        \mu(i) \saddR(i) = 0 \cdot
    \end{gather}

    By summing Equation~\ref{eq:kkt_proof_R_is_pxIe} over $x$ and using Equation~\ref{eq:kkt_proof_r_sums_one},
    we have
    \[
        \eta = 1 \cdot
    \]
    Thus, we get Equation~\ref{eq:kkt_w_avg} which also implicitly implies Equation~\ref{eq:kkt_proof_p_sums_one}.
\end{proof}

\begin{theorem} [Centrality]
    \label{th:lambda-centrality}
    Let $X$ be a discrete random variable, and let \setQ be a finite set of distributions over $X$.
    Suppose \cDist is a convex combination of elements of \setQ, with all weights strictly positive.
    If for every \inSet{Q} and some constant $\lambda$ we have:
    \begin{equation*}
        \kl[][\cDist] = \lambda,
    \end{equation*}
    then \cDist is a central distribution of \setQ for $X$.
\end{theorem}

\begin{proof}
    Lemma~\ref{lm:saddle_conditions} provides a set of sufficient conditions for a saddle point.
    Note that more restrictive conditions can also be imposed to derive an alternative set of sufficient conditions.
    We can replace Inequality~\ref{eq:saddle_lm_r_geq_0} with the more
    restrictive $\saddR(i) > 0$ condition.
    By Equation~\ref{eq:saddle_lm_rmu_is_zero}, that implies $\mu(i) = 0$,
    and Lemma~\ref{lm:saddle_conditions} conditions can be simplified to
    \begin{gather*}
        \klExp[Q][\saddP][x] = \lambda, \\
        \saddP(x) = \sum_i \saddR(i) Q(x),
    \end{gather*}
    which must hold for every \inSet{Q} and a positive probability distribution \saddR.

    Since \cDist satisfies these condition, Lemma~\ref{lm:saddle_conditions} and
    Theorem~\ref{th:saddle_is_central} imply that \cDist is a central distribution of \setQ for $X$.
\end{proof}

These theorems were proved under the assumption that \setQ is finite; however, they can be extended to cases where
\setQ can be indexed by a finite-dimensional real vector \param.
We can treat this vector as a random variable and interpret $Q_{\paramVal}$ as a conditional
distribution $Q(\xIt)$.
Then, the weight function, $\saddR(i)$, will essentially become a prior distribution for \paramVal.

\begin{theorem}
    Let $X$ be a discrete random variable and \setQ be a set of distributions over $X$, indexed by a
    finite-dimensional random vector \param.
    Let \cPrior be a positive prior distribution for \param such that for $\like(\xIt) \in \setQ$ we have
    \begin{equation*}
        \cPrior(\paramVal) \propto \exp \left\{  \sum_{x} \like(\xIt) \ln \cPrior(\tIx) \right\} \cdot
    \end{equation*}
    Then, $\cDist(x) = \mrg$ is a central distribution of \setQ for $X$.
    We also refer to \cPrior as a central prior for $X$.
\end{theorem}

\begin{proof}
    To Do!!!
\end{proof}